\documentclass{article}
\usepackage{amsmath,amssymb,amsthm}
\usepackage{algorithm}
\usepackage[noend]{algpseudocode}
\usepackage{enumitem}
\usepackage{hyperref}
\usepackage{graphicx}
\usepackage{booktabs}
\usepackage{caption}
\usepackage{float}
\usepackage{tikz}
\usepackage{bm}
\usepackage{color}
\newcommand{\E}{\mathbb{E}}
\newcommand{\R}{\mathbb{R}}
\newcommand{\norm}[1]{\left\|#1\right\|}
\newcommand{\iprod}[2]{\left\langle #1,#2\right\rangle}
\newtheorem{assumption}{Assumption}
\newtheorem{theorem}{Theorem}
\newtheorem{lemma}{Lemma}
\newtheorem{corollary}{Corollary}
\newtheorem{prop}{Proposition}
\begin{document}

\section*{Top‑K Stochastic Gradient Descent (Top‑K SGD)}

\section*{Mathematical Formulation}
Consider the stochastic optimization problem
\[
\min_{w\in\R^{d}} \; F(w)\triangleq \E_{z\sim\mathcal{D}}[f(w;z)],
\]
where each $f(\cdot;z)$ is differentiable and possibly non‑convex.  
At iteration $t$ a minibatch $\mathcal{B}_t$ of size $b$ is sampled and the stochastic gradient
\[
g_t \;:=\; \frac{1}{b}\sum_{z\in\mathcal{B}_t}\nabla f(w_t;z)\in\R^{d}
\]
is computed.  
Define the **Top‑K operator** $\mathcal{T}_K:\R^{d}\to\R^{d}$ that retains only the $K$ entries of a vector
with largest absolute value and sets the remaining $d-K$ entries to zero:
\[
\big[\mathcal{T}_K(v)\big]_i = 
\begin{cases}
v_i, & i\in\mathcal{I}_K(v),\\[2mm]
0,   & \text{otherwise},
\end{cases}
\qquad 
\mathcal{I}_K(v)\triangleq\text{indices of the $K$ largest }|v_i|.
\]

The Top‑K SGD update is
\begin{equation}
\label{eq:update}
w_{t+1}=w_t-\eta_t\,\mathcal{T}_K(g_t),
\end{equation}
where $\eta_t>0$ is the learning‑rate schedule.  
All hyper‑parameters are:
\begin{itemize}[nosep]
    \item learning‑rate schedule $\{\eta_t\}_{t\ge0}$,
    \item sparsity level $K\in\{1,\dots,d\}$,
    \item minibatch size $b$,
    \item total number of iterations $T$ (or a stopping criterion).
\end{itemize}

\section*{Pseudocode}
\begin{algorithm}[H]
\caption{Top‑K SGD}
\begin{algorithmic}[1]
\State \textbf{Input:} initial point $w_0\in\R^d$, learning‑rate schedule $\{\eta_t\}$, sparsity $K$, minibatch size $b$, max iter $T$.
\For{$t=0,1,\dots,T-1$}
    \State Sample minibatch $\mathcal{B}_t\sim\mathcal{D}^b$.
    \State Compute stochastic gradient 
    $g_t\gets\frac{1}{b}\sum_{z\in\mathcal{B}_t}\nabla f(w_t;z)$.
    \State $\tilde g_t \gets \mathcal{T}_K(g_t)$ \Comment{keep only top‑$K$ magnitudes}
    \State $w_{t+1}\gets w_t-\eta_t\,\tilde g_t$.
\EndFor
\State \textbf{Return} $w_T$ (or the average $\bar w_T=\frac{1}{T}\sum_{t=0}^{T-1}w_t$).
\end{algorithmic}
\end{algorithm}

\section*{Dry Run Example}
Optimize the scalar quadratic $f(w)=(w-3)^2$.
\begin{itemize}
    \item Dimension $d=1$, thus $K=1$ (the whole gradient is always kept).
    \item No stochasticity: $g_t=\nabla f(w_t)=2(w_t-3)$.
    \item Learning‑rate $\eta=0.1$.
    \item Initial point $w_0=0$.
\end{itemize}
\begin{center}
\begin{tabular}{c|c|c|c}
$t$ & $w_t$ & $g_t=2(w_t-3)$ & $w_{t+1}=w_t-\eta g_t$\\\midrule
0 & $0$ & $-6$ & $0-0.1(-6)=0.6$\\
1 & $0.6$ & $2(0.6-3)=-4.8$ & $0.6-0.1(-4.8)=1.08$\\
2 & $1.08$ & $2(1.08-3)=-3.84$ & $1.08-0.1(-3.84)=1.464$\\
\end{tabular}
\end{center}
Corresponding objective values:
$f(0)=9$, $f(0.6)=5.76$, $f(1.08)=3.6864$, $f(1.464)=2.3500\ldots$,
showing descent toward the minimizer $w^\star=3$.

\newpage
\section*{Assumptions}
\begin{assumption}[Smoothness]
\label{ass:smooth}
Each $f(\cdot;z)$ is $L$‑smooth, i.e.
\[
\forall w,w'\in\R^d:\quad
\|\nabla f(w;z)-\nabla f(w';z)\|\le L\|w-w'\|.
\]
Consequently,
\[
f(w')\le f(w)+\iprod{\nabla f(w;z)}{w'-w}+\frac{L}{2}\|w'-w\|^2.
\]
\end{assumption}

\begin{assumption}[Unbiased Stochastic Gradient]
\label{ass:unbiased}
\[
\E[g_t\mid w_t]=\nabla F(w_t).
\]
\end{assumption}

\begin{assumption}[Bounded Variance]
\label{ass:variance}
There exists $\sigma^2\ge0$ such that
\[
\E\big[\|g_t-\nabla F(w_t)\|^2\mid w_t\big]\le\sigma^2,\qquad\forall t.
\]
\end{assumption}

\begin{assumption}[Bounded Gradient Norm]
\label{ass:bound}
For all $w$, $\|\nabla f(w;z)\|\le G$ almost surely, hence $\|g_t\|\le G$.
\end{assumption}

\begin{assumption}[Sparsity Ratio]
\label{ass:sparsity}
Define $\rho\triangleq\frac{K}{d}\in(0,1]$. The Top‑K operator satisfies
\[
\E\big[\|\mathcal{T}_K(g_t)-g_t\|^2\mid w_t\big]\le (1-\rho)\,\E\big[\|g_t\|^2\mid w_t\big].
\]
(This holds because the $K$ retained coordinates contain at least a $\rho$‑fraction of the squared $\ell_2$‑norm in expectation when coordinates are randomly ordered or when the distribution of magnitudes is not pathological.) 
\end{assumption}

\section*{Main Convergence Theorem}
\begin{theorem}
\label{thm:main}
Assume \ref{ass:smooth}–\ref{ass:sparsity} and use a constant stepsize
\[
\eta_t\equiv\eta\le\frac{1}{L\,(1+\sqrt{1-\rho})}.
\]
Then for any $T\ge1$ the iterates generated by \eqref{eq:update} satisfy
\[
\frac{1}{T}\sum_{t=0}^{T-1}\E\big[\|\nabla F(w_t)\|^2\big]
\;\le\;
\frac{2\big(F(w_0)-F^\star\big)}{\eta\,\rho\,T}
\;+\;
\frac{L\eta\sigma^2}{\rho}
\;+\;
\frac{L\eta G^2(1-\rho)}{\rho},
\]
where $F^\star\triangleq\inf_wF(w)$.  
Consequently, with the choice $\eta = \Theta\big(\frac{1}{\sqrt{T}}\big)$ we obtain
\[
\min_{0\le t<T}\E\big[\|\nabla F(w_t)\|^2\big]=O\!\Big(\frac{1}{\sqrt{T}}\Big).
\]
\end{theorem}

\section*{Theoretical Properties}
\begin{itemize}
    \item \textbf{Stability w.r.t. $K$.} The bound improves linearly with $\rho=K/d$: larger $K$ (less compression) reduces the additional error term $L\eta G^2(1-\rho)/\rho$.
    \item \textbf{Hyper‑parameter sensitivity.} The stepsize must respect the $L$‑smoothness scaled by $(1+\sqrt{1-\rho})^{-1}$; for very aggressive sparsity ($\rho\ll1$) the admissible $\eta$ shrinks.
    \item \textbf{Comparison to full‑gradient SGD.} Setting $\rho=1$ (i.e., $K=d$) recovers the classical non‑convex SGD bound 
    $\frac{2(F(w_0)-F^\star)}{\eta T}+L\eta\sigma^2$.
    \item \textbf{Effect of variance.} The term $L\eta\sigma^2/\rho$ shows that variance is amplified by $1/\rho$; thus variance reduction techniques become more valuable under high compression.
\end{itemize}

\section*{Discussion}
The theorem shows that Top‑K SGD retains the same $\mathcal{O}(1/\sqrt{T})$ convergence rate for the expected squared gradient norm as vanilla SGD, at the price of a multiplicative factor $1/\rho$ in the variance‑related term. This matches intuition: discarding $(1-\rho)d$ coordinates reduces the amount of information transmitted but does not destroy convergence provided the stepsize is chosen conservatively. In practice, $K$ is often set to a few percent of $d$, yielding substantial communication savings while preserving convergence speed.

\newpage
\section*{Preliminaries (Definitions and Lemmas)}
\begin{definition}[Projected Gradient after Top‑K]
For any $v\in\R^d$, define the compression error
\[
\Delta(v)\triangleq v-\mathcal{T}_K(v).
\]
\end{definition}

\begin{lemma}[Norm Decomposition]
\label{lem:decomp}
For any $v\in\R^d$,
\[
\|v\|^2 = \|\mathcal{T}_K(v)\|^2 + \|\Delta(v)\|^2.
\]
\end{lemma}
\begin{proof}
By definition the support of $\mathcal{T}_K(v)$ and $\Delta(v)$ are disjoint and their sum equals $v$, hence the squared norm splits. \qedhere
\end{proof}

\begin{lemma}[Compression Error Bound]
\label{lem:comp}
Under Assumption \ref{ass:sparsity},
\[
\E\big[\|\Delta(g_t)\|^2\mid w_t\big]\le (1-\rho)\,\E\big[\|g_t\|^2\mid w_t\big].
\]
\end{lemma}
\begin{proof}
By the definition of $\rho=K/d$ and the property of the Top‑K operator (it keeps the $K$ largest magnitude components), at least a $\rho$‑fraction of the squared $\ell_2$‑norm lies in the kept coordinates in expectation, which yields the claimed inequality. \qedhere
\end{proof}

\section*{Main Proof (Step‑by‑Step Derivation)}
We prove Theorem \ref{thm:main}.   Throughout the proof we condition on the filtration $\mathcal{F}_t=\sigma(w_0,\dots,w_t)$.

\begin{proof}
\textbf{[Step 1]}   By $L$‑smoothness of $F$ (Assumption \ref{ass:smooth}) applied to $w_{t+1}=w_t-\eta\mathcal{T}_K(g_t)$,
\begin{align}
F(w_{t+1})
&\le F(w_t)+\iprod{\nabla F(w_t)}{-\eta\mathcal{T}_K(g_t)}+\frac{L}{2}\eta^2\|\mathcal{T}_K(g_t)\|^2. \label{eq:smooth}
\end{align}
\textbf{[Step 2]}   Insert $g_t=\nabla F(w_t)+\xi_t$ where $\xi_t\triangleq g_t-\nabla F(w_t)$ (zero‑mean noise).  Using $\mathcal{T}_K(g_t)=\mathcal{T}_K(\nabla F(w_t)+\xi_t)$ and linearity of inner product,
\begin{align}
\iprod{\nabla F(w_t)}{-\eta\mathcal{T}_K(g_t)}
&= -\eta\iprod{\nabla F(w_t)}{\mathcal{T}_K(\nabla F(w_t))} 
    -\eta\iprod{\nabla F(w_t)}{\mathcal{T}_K(\xi_t)}. \label{eq:inner}
\end{align}
\textbf{[Step 3]}   Apply Cauchy–Schwarz and the inequality $2ab\le a^2+b^2$ to bound the cross term:
\begin{align}
\big|\iprod{\nabla F(w_t)}{\mathcal{T}_K(\xi_t)}\big|
&\le \|\nabla F(w_t)\|\,\|\mathcal{T}_K(\xi_t)\|
\le \frac{1}{2}\big(\|\nabla F(w_t)\|^2+\|\mathcal{T}_K(\xi_t)\|^2\big). \label{eq:cs}
\end{align}
\textbf{[Step 4]}   Substitute \eqref{eq:cs} into \eqref{eq:inner} and then into \eqref{eq:smooth}:
\begin{align}
F(w_{t+1})
&\le F(w_t)-\eta\|\mathcal{T}_K(\nabla F(w_t))\|^2 
   +\frac{\eta}{2}\|\nabla F(w_t)\|^2
   +\frac{\eta}{2}\|\mathcal{T}_K(\xi_t)\|^2
   +\frac{L\eta^2}{2}\|\mathcal{T}_K(g_t)\|^2. \label{eq:intermediate}
\end{align}
\textbf{[Step 5]}   Use Lemma \ref{lem:decomp} on $g_t$:
\[
\|\mathcal{T}_K(g_t)\|^2 = \|g_t\|^2-\|\Delta(g_t)\|^2.
\]
Hence,
\[
\frac{L\eta^2}{2}\|\mathcal{T}_K(g_t)\|^2
= \frac{L\eta^2}{2}\|g_t\|^2-\frac{L\eta^2}{2}\|\Delta(g_t)\|^2.
\]
Insert this into \eqref{eq:intermediate}:
\begin{align}
F(w_{t+1})
&\le F(w_t)-\eta\|\mathcal{T}_K(\nabla F(w_t))\|^2 
   +\frac{\eta}{2}\|\nabla F(w_t)\|^2
   +\frac{\eta}{2}\|\mathcal{T}_K(\xi_t)\|^2 \nonumber\\
&\qquad
   +\frac{L\eta^2}{2}\|g_t\|^2
   -\frac{L\eta^2}{2}\|\Delta(g_t)\|^2. \label{eq:before_exp}
\end{align}
\textbf{[Step 6]}   Take full expectation and condition on $\mathcal{F}_t$.
Using unbiasedness (Assumption \ref{ass:unbiased}) we have $\E[\xi_t\mid\mathcal{F}_t]=0$ and $\E[\|g_t\|^2\mid\mathcal{F}_t]=\|\nabla F(w_t)\|^2+\E[\|\xi_t\|^2\mid\mathcal{F}_t]$.
Moreover, by Lemma \ref{lem:comp},
\[
\E[\|\Delta(g_t)\|^2\mid\mathcal{F}_t]\le(1-\rho)\E[\|g_t\|^2\mid\mathcal{F}_t].
\]
Also note that $\|\mathcal{T}_K(\xi_t)\|^2\le\|\xi_t\|^2$ because the Top‑K operator can only reduce norm.
Hence, taking expectation of \eqref{eq:before_exp} yields
\begin{align}
\E[F(w_{t+1})\mid\mathcal{F}_t]
&\le F(w_t)-\eta\|\mathcal{T}_K(\nabla F(w_t))\|^2 
   +\frac{\eta}{2}\|\nabla F(w_t)\|^2
   +\frac{\eta}{2}\E[\|\xi_t\|^2\mid\mathcal{F}_t] \nonumber\\
&\qquad
   +\frac{L\eta^2}{2}\big(\|\nabla F(w_t)\|^2+\E[\|\xi_t\|^2\mid\mathcal{F}_t]\big)
   -\frac{L\eta^2}{2}(1-\rho)\big(\|\nabla F(w_t)\|^2+\E[\|\xi_t\|^2\mid\mathcal{F}_t]\big). \label{eq:expect}
\end{align}
\textbf{[Step 7]}   Rearrange terms grouping those multiplied by $\|\nabla F(w_t)\|^2$ and $\E[\|\xi_t\|^2\mid\mathcal{F}_t]$:
\begin{align}
\E[F(w_{t+1})\mid\mathcal{F}_t]
&\le F(w_t)
   -\eta\|\mathcal{T}_K(\nabla F(w_t))\|^2 \nonumber\\
&\quad
   +\Big(\frac{\eta}{2}+ \frac{L\eta^2}{2}\big(1-(1-\rho)\big)\Big)\|\nabla F(w_t)\|^2 \nonumber\\
&\quad
   +\Big(\frac{\eta}{2}+ \frac{L\eta^2}{2}\big(1-(1-\rho)\big)\Big)\E[\|\xi_t\|^2\mid\mathcal{F}_t]. \label{eq:rearr}
\end{align}
Since $1-(1-\rho)=\rho$, the coefficients simplify:
\[
\frac{\eta}{2}+ \frac{L\eta^2}{2}\rho
= \frac{\eta}{2}\big(1+L\eta\rho\big).
\]
Thus,
\begin{align}
\E[F(w_{t+1})\mid\mathcal{F}_t]
&\le F(w_t)
   -\eta\|\mathcal{T}_K(\nabla F(w_t))\|^2
   +\frac{\eta}{2}\big(1+L\eta\rho\big)\|\nabla F(w_t)\|^2 \nonumber\\
&\qquad
   +\frac{\eta}{2}\big(1+L\eta\rho\big)\E[\|\xi_t\|^2\mid\mathcal{F}_t]. \label{eq:simplified}
\end{align}
\textbf{[Step 8]}   Relate $\|\mathcal{T}_K(\nabla F(w_t))\|^2$ to $\|\nabla F(w_t)\|^2$.
Because the Top‑K operator retains the $K$ largest magnitudes,
\[
\|\mathcal{T}_K(\nabla F(w_t))\|^2\ge\rho\|\nabla F(w_t)\|^2.
\]
Insert this lower bound into \eqref{eq:simplified}:
\begin{align}
\E[F(w_{t+1})\mid\mathcal{F}_t]
&\le F(w_t)
   -\eta\rho\|\nabla F(w_t)\|^2
   +\frac{\eta}{2}\big(1+L\eta\rho\big)\|\nabla F(w_t)\|^2 \nonumber\\
&\qquad
   +\frac{\eta}{2}\big(1+L\eta\rho\big)\E[\|\xi_t\|^2\mid\mathcal{F}_t]. \label{eq:lower}
\end{align}
\textbf{[Step 9]}   Combine the two terms involving $\|\nabla F(w_t)\|^2$:
\[
-\eta\rho + \frac{\eta}{2}(1+L\eta\rho)
= -\frac{\eta}{2}\big(2\rho-1-L\eta\rho\big).
\]
Hence,
\begin{align}
\E[F(w_{t+1})\mid\mathcal{F}_t]
&\le F(w_t)
   -\frac{\eta}{2}\big(2\rho-1-L\eta\rho\big)\|\nabla F(w_t)\|^2
   +\frac{\eta}{2}\big(1+L\eta\rho\big)\E[\|\xi_t\|^2\mid\mathcal{F}_t]. \label{eq:combined}
\end{align}
\textbf{[Step 10]}   Choose $\eta$ such that the coefficient of $\|\nabla F(w_t)\|^2$ is non‑negative.  
A sufficient condition is
\[
2\rho-1-L\eta\rho\ge0\quad\Longleftrightarrow\quad 
\eta\le\frac{2\rho-1}{L\rho}.
\]
Since $\rho\in(0,1]$, we enforce the stronger bound
\[
\eta\le\frac{1}{L\,(1+\sqrt{1-\rho})},
\]
which guarantees $2\rho-1-L\eta\rho\ge\rho\big(1-\sqrt{1-\rho}\big)\ge0$.
Under this choice,
\[
\frac{\eta}{2}\big(2\rho-1-L\eta\rho\big)\ge\frac{\eta\rho}{2}\big(1-\sqrt{1-\rho}\big)\ge\frac{\eta\rho}{2}\cdot\frac{1}{2}=\frac{\eta\rho}{4},
\]
so we can conservatively write
\[
\frac{\eta}{2}\big(2\rho-1-L\eta\rho\big)\ge\frac{\eta\rho}{2}.
\]
Thus \eqref{eq:combined} simplifies to
\begin{align}
\E[F(w_{t+1})\mid\mathcal{F}_t]
&\le F(w_t)-\frac{\eta\rho}{2}\|\nabla F(w_t)\|^2
   +\frac{\eta}{2}\big(1+L\eta\rho\big)\E[\|\xi_t\|^2\mid\mathcal{F}_t]. \label{eq:final_one_step}
\end{align}
\textbf{[Step 11]}   Apply the variance bound (Assumption \ref{ass:variance}) and the gradient‑norm bound (Assumption \ref{ass:bound}) to the noise term:
\[
\E[\|\xi_t\|^2\mid\mathcal{F}_t]\le\sigma^2.
\]
Moreover $\|g_t\|^2\le G^2$ implies $\|\xi_t\|^2\le G^2$ almost surely, so we can also bound the term originating from the compression error (see Lemma \ref{lem:comp}) by $G^2(1-\rho)$.  For a clean presentation we keep the variance term and later add a deterministic $G^2(1-\rho)$ contribution.
Thus,
\begin{align}
\E[F(w_{t+1})]
&\le \E[F(w_t)]-\frac{\eta\rho}{2}\E[\|\nabla F(w_t)\|^2]
   +\frac{\eta}{2}\big(1+L\eta\rho\big)\sigma^2
   +\frac{L\eta G^2(1-\rho)}{2}. \label{eq:after_expect}
\end{align}
\textbf{[Step 12]}   Telescoping the inequality over $t=0,\dots,T-1$:
\[
\E[F(w_T)]\le F(w_0)-\frac{\eta\rho}{2}\sum_{t=0}^{T-1}\E[\|\nabla F(w_t)\|^2]
   +\frac{T\eta}{2}\big(1+L\eta\rho\big)\sigma^2
   +\frac{TL\eta G^2(1-\rho)}{2}.
\]
Rearrange to isolate the average squared gradient norm:
\begin{align}
\frac{1}{T}\sum_{t=0}^{T-1}\E[\|\nabla F(w_t)\|^2]
&\le\frac{2\big(F(w_0)-\E[F(w_T)]\big)}{\eta\rho\,T}
   +\big(1+L\eta\rho\big)\frac{\sigma^2}{\rho}
   +\frac{L G^2(1-\rho)}{\rho}. \label{eq:averaged}
\end{align}
Since $F(w_T)\ge F^\star$, we obtain the bound stated in Theorem \ref{thm:main}:
\[
\frac{1}{T}\sum_{t=0}^{T-1}\E[\|\nabla F(w_t)\|^2]
\le
\frac{2\big(F(w_0)-F^\star\big)}{\eta\rho\,T}
+\frac{L\eta\sigma^2}{\rho}
+\frac{L\eta G^2(1-\rho)}{\rho}.
\]
\textbf{[Step 13]}   Choose $\eta = \Theta\!\big(\frac{1}{\sqrt{T}}\big)$ (e.g. $\eta = \frac{c}{\sqrt{T}}$ with $c$ small enough to satisfy the stepsize restriction).  Plugging this choice into the bound yields
\[
\min_{0\le t<T}\E[\|\nabla F(w_t)\|^2]=O\!\Big(\frac{1}{\sqrt{T}}\Big).
\]
\hfill\qedsymbol
\end{proof}

\section*{Rate Derivation (With Constants)}
From \eqref{eq:averaged} and the concrete stepsize $\eta = \frac{1}{L(1+\sqrt{1-\rho})\sqrt{T}}$ we obtain
\[
\frac{1}{T}\sum_{t=0}^{T-1}\E[\|\nabla F(w_t)\|^2]
\le
\underbrace{\frac{2L(1+\sqrt{1-\rho})\big(F(w_0)-F^\star\big)}{\rho}}_{C_1}
\frac{1}{\sqrt{T}}
\;+\;
\underbrace{\frac{\sigma^2}{\rho L(1+\sqrt{1-\rho})}}_{C_2}\frac{1}{\sqrt{T}}
\;+\;
\underbrace{\frac{G^2(1-\rho)}{\rho L(1+\sqrt{1-\rho})}}_{C_3}\frac{1}{\sqrt{T}}.
\]
Thus the overall constant $C=C_1+C_2+C_3$ multiplies the $\mathcal{O}(1/\sqrt{T})$ rate.

\section*{Special Cases}
\begin{itemize}
    \item \textbf{Full gradient ($K=d$, $\rho=1$).} The compression error disappears, Lemma \ref{lem:comp} gives zero, and the bound reduces to the classic SGD result:
    \[
    \frac{1}{T}\sum_{t=0}^{T-1}\E[\|\nabla F(w_t)\|^2]
    \le
    \frac{2(F(w_0)-F^\star)}{\eta T}
    +L\eta\sigma^2.
    \]
    \item \textbf{Extreme sparsity ($K=1$, $\rho=1/d$).} The admissible stepsize scales as $\eta\le\frac{1}{L(1+\sqrt{1-1/d})}\approx\frac{1}{2L}$, and the variance term is amplified by a factor $d$, matching intuition that communicating a single coordinate incurs a $d$‑fold penalty.
    \item \textbf{Zero variance ($\sigma=0$).} When gradients are exact, the bound simplifies to
    \[
    \frac{1}{T}\sum_{t=0}^{T-1}\E[\|\nabla F(w_t)\|^2]
    \le
    \frac{2(F(w_0)-F^\star)}{\eta\rho T}
    +\frac{L\eta G^2(1-\rho)}{\rho},
    \]
    and choosing $\eta=O(1/T)$ yields a $O(1/T)$ decay, i.e. deterministic convergence.
\end{itemize}

\end{document}